\section*{Capítulo \ref{ch:b1:magnetostatics} --- Magnetostática: campo, fuerzas y dipolos}

\begin{solution}{exo:b1:magnetostatics:1}
$B = \mu_0I/2\pi r = 2 \times 10^{-7} \times 10/0.01 = \qty{2e-4}{T}$. Para \qty{1}{T}:
$I = 2\pi rB/\mu_0 = \qty{5e4}{A}$. El valor terrestre, a $r = 2 \times 10^{-7} \times
100/5 \times 10^{-5} = \qty{0.4}{m}$.
\end{solution}

\begin{solution}{exo:b1:magnetostatics:2}
Centro: $\mu_0NI/2R = 4\pi \times 10^{-7} \times 200/0.1 = \qty{2.5}{mT}$. En $z = R$:
$\times(R^2/2R^2)^{3/2} = 2^{-3/2}$: \qty{0.89}{mT}. En $z = 10R$: $\times(1/101)^{3/2}
= 10^{-3}$: \qty{2.5}{\micro T}, ya la ley dipolar en $1/z^3$.
\end{solution}

\begin{solution}{exo:b1:magnetostatics:3}
$B = \mu_0nI = 4\pi \times 10^{-7} \times 1000 \times 5 = \qty{6.3}{mT}$; en el extremo,
la mitad: \qty{3.1}{mT}. Con hilo de \qty{0.5}{mm}: $2000$ espiras por metro, luego
\qty{12.6}{mT} con \qty{5}{A}.
\end{solution}

\begin{solution}{exo:b1:magnetostatics:4}
$F = ILB = 10 \times 1 \times 0.1 = \qty{1.0}{N}$; a \qty{30}{\degree}, $\sin30^\circ$:
\qty{0.5}{N}. Motor: $100 \times \qty{1}{N} \times \qty{0.05}{m} = \qty{5}{N.m}$ ---
un motor industrial pequeño.
\end{solution}

\begin{solution}{exo:b1:magnetostatics:5}
Con $z = d\tan\alpha$: $B = \dfrac{\mu_0I}{4\pi}\displaystyle\int\dfrac{d\,\dd z}{(d^2 + z^2)^{3/2}}
= \dfrac{\mu_0I}{4\pi d}\displaystyle\int_{\alpha_1}^{\alpha_2}\cos\alpha\,\dd\alpha = \dfrac{\mu_0I}{4\pi d}
(\sin\alpha_2 - \sin\alpha_1)$; hilo infinito, $\alpha_{1,2} = \mp\pi/2$: $\mu_0I/2\pi d$.
Cuadrado: cada lado a $d = a/2$ visto bajo $\pm45^\circ$ da $\mu_0I\sqrt2/2\pi a$;
los cuatro lados: $B = 2\sqrt2\,\mu_0I/\pi a = 0.90\,\mu_0I/a$ (una circunferencia del mismo
``radio'' $a/2$ daría $\mu_0I/a$).
\end{solution}

\begin{solution}{exo:b1:magnetostatics:6}
$B(z) = f(z - R/2) + f(z + R/2)$ con $f(u) = \mu_0NIR^2/2(R^2 + u^2)^{3/2}$.
Punto medio: $2f(R/2) = \mu_0NIR^2/(5R^2/4)^{3/2} = (4/5)^{3/2}\mu_0NI/R$. $f'$ es
impar, luego $B'(0) = f'(-R/2) + f'(R/2) = 0$; $f''(u) \propto (4u^2 - R^2)(R^2 +
u^2)^{-7/2}$ se anula en $u = R/2$, luego $B''(0) = 2f''(R/2) = 0$: el campo es
uniforme hasta el tercer orden. Cifras: $0.716 \times 4\pi \times 10^{-7} \times 100/0.15 =
\qty{0.60}{mT}$.
\end{solution}

\begin{solution}{exo:b1:magnetostatics:7}
Superficie: $\mu_0I/2\pi a = 2 \times 10^{-7} \times 100/0.002 = \qty{10}{mT}$; a
\qty{1}{mm}, dentro, $\times r/a$: \qty{5}{mT}; a \qty{1}{cm}: \qty{2}{mT}. $B$ crece
linealmente hasta la superficie y luego cae como $1/r$. Coaxial: fuera de la malla
la corriente enlazada es nula, $B = 0$.
\end{solution}

\begin{solution}{exo:b1:magnetostatics:8}
$2 \times 10^{-7}$~N/m. Barras: $2 \times 10^{-7} \times 10^8/0.1 = \qty{200}{N/m}$;
a \qty{100}{kA}: $\qty{2e4}{N/m}$ --- dos toneladas por metro. Los soportes
aislantes de la aparamenta se dimensionan para la corriente de defecto, no
para la de servicio.
\end{solution}

\begin{solution}{exo:b1:magnetostatics:9}
$\Gamma = mB_h = \qty{2e-6}{N.m}$; $W = 2mB_h = \qty{4e-6}{J}$; $J\ddot\theta =
-mB_h\theta$: $T = 2\pi\sqrt{J/mB_h} = 2\pi\sqrt{10^{-7}/2 \times 10^{-6}} =
\qty{1.4}{s}$. Se cronometran las oscilaciones: $B_h = 4\pi^2J/mT^2$ una vez
conocidos $m$ y $J$ (Gauss medía $m$ aparte, por la desviación que produce
en una segunda aguja).
\end{solution}

\begin{solution}{exo:b1:magnetostatics:10}
(a) $\mu_0IR^2/2z^3 = \mu_0(I\pi R^2)/2\pi z^3 = \mu_0m/2\pi z^3 = (\mu_0/4\pi)\,2m/z^3$:
el eléctrico $2p/4\pi\varepsilon_0z^3$ con $1/\varepsilon_0 \to \mu_0$. (b) $m = 2\pi R_E^3B/
\mu_0 = 2\pi \times 2.58 \times 10^{20} \times 6 \times 10^{-5}/1.26 \times 10^{-6} = \qty{7.7e22}{A.m^2}$.
(c) $I = m/\pi R^2 = 7.7 \times 10^{22}/(\pi \times 9 \times 10^{12}) = \qty{2.7e9}{A}$. (d)
Ecuador ($\theta = 90^\circ$): $B_\theta = \mu_0m/4\pi R_E^3 = \qty{3e-5}{T}$, la mitad del
valor polar. En $\theta = 45^\circ$: $\tan(\text{incl}) = B_r/B_\theta = 2\cot\theta = 2$,
inclinación $= \qty{63}{\degree}$ --- el campo se hunde con fuerte pendiente en el
suelo en las latitudes medias.
\end{solution}

\begin{solution}{exo:b1:magnetostatics:11}
$\mu_0nI = 4\pi \times 10^{-7} \times 2000 \times 3 = \qty{7.5}{mT}$. Centro: $\tan\alpha_1 =
R/(L/2) = 0.2$, $\cos\alpha_1 = -\cos\alpha_2 = 0.981$: $B = 0.981\,\mu_0nI =
\qty{7.4}{mT}$ (el $98\%$). Extremo: $\cos\alpha_1 = 0$, $\tan\alpha_2 = R/L$, $\cos\alpha_2 =
-0.995$: $B = 0.497\,\mu_0nI = \qty{3.7}{mT}$ (el $50\%$). Fuera, a una distancia $x$ del
extremo, con los dos extremos del mismo lado: $B = \tfrac12\mu_0nI(\cos\alpha_1 - \cos\alpha_2)
\approx \tfrac12\mu_0nI\,\tfrac{R^2}{2}[1/x^2 - 1/(x + L)^2]$; el $1\%$ del valor
central en $x \approx \qty{10}{cm}$, media longitud fuera. $B(z)$: una meseta plana en
casi toda la longitud, que baja a la mitad en los extremos y a unos pocos
por ciento a una longitud de distancia.
\end{solution}

\begin{solution}{exo:b1:magnetostatics:12}
(a) El campo propio del devanado salta de $0$ a $B$ al atravesarlo; el
campo que actúa sobre él es la media, $B/2$ (la capa no puede empujarse a
sí misma). Fuerza sobre un elemento de devanado de longitud $\dd l$ y anchura $\dd z$
que lleva $K\,\dd z$: $K\,\dd z\,\dd l\,B/2$, radial hacia fuera (compruébese con $I\,\dd\vect l
\wedge\vect B$): presión $p = KB/2 = B^2/2\mu_0$, ya que $K = nI = B/\mu_0$.
(b) \qty{1}{T}: \qty{4e5}{Pa} = \qty{4}{bar}; \qty{10}{T}: \qty{400}{bar}; \qty{45}{T}:
\qty{8e8}{Pa} = \qty{8000}{bar}, cerca del \qty{1}{GPa} del acero. (c) El
devanado ha de soportar una presión que crece como $B^2$: por encima de \qty{40}{T}
más o menos, ningún material ni sistema de refrigeración aguanta --- el límite de
los imanes permanentes es mecánico (y, en los superconductores, su campo crítico).
\end{solution}

\begin{solution}{pb:b1:magnetostatics:1}
\textbf{1.} Todo plano que contiene al eje es un \omterm{prop:b1:electrostatics-gauss:symmetry}{plano de simetría} de las
corrientes: $\vect B$ es perpendicular a él, ortorradial; por invariancia
a lo largo y alrededor del eje: $\vect B = B(r)\,\vect e_\theta$.

\textbf{2.} Circunferencia de radio $r$: $2\pi rB = \mu_0I_{\text{enc}}$. Alma:
$I_{\text{enc}} = Ir^2/a^2$, $B = \mu_0Ir/2\pi a^2$; entre ambos: $\mu_0I/2\pi r$;
fuera ($r > c$): $I - I = 0$, $B = 0$.

\textbf{3.} $B(a) = 2 \times 10^{-7}/5 \times 10^{-4} = \qty{4.0e-4}{T}$; $B(b) = 2 \times
10^{-7}/2.5 \times 10^{-3} = \qty{8.0e-5}{T}$.

\textbf{4.} Dentro del tubo no hay corriente enlazada: allí $B = 0$ en vez
del crecimiento lineal; en todo lo demás nada cambia --- el campo fuera de
una corriente cilíndrica solo depende de la corriente total.

\textbf{5.} $I_{\text{enc}} = I - I(r^2 - b^2)/(c^2 - b^2) = I(c^2 - r^2)/(c^2 -
b^2)$: $B = \dfrac{\mu_0I}{2\pi r}\,\dfrac{c^2 - r^2}{c^2 - b^2}$, igual a $\mu_0I/2\pi b$ en
$b$ y a $0$ en $c$: continuo. Esbozo: lineal hasta \qty{4e-4}{T} en
\qty{0.5}{mm}, en $1/r$ hasta \qty{8e-5}{T} en \qty{2.5}{mm}, cayendo a cero en
\qty{3}{mm} y nulo más allá.

\textbf{6.} La corriente de la malla es antiparalela a la del alma: repulsión,
una presión hacia fuera. Corriente superficial $K = I/2\pi b = \qty{64}{A/m}$; campo medio
$\tfrac12(8 \times 10^{-5} + 0) = \qty{4e-5}{T}$: $p = KB = \qty{2.5e-3}{Pa}$. A
\qty{10}{kA}: $\times10^8$, \qty{2.5e5}{Pa} = \qty{2.5}{bar} --- los cables de
potencia pulsada se construyen para aguantarlo.

\textbf{7.} No hay campo fuera: el cable ni radia ni recoge interferencias
magnéticas de sus vecinos. Un par de hilos: los dos campos casi se
cancelan y dejan $\approx \mu_0Id/2\pi D^2 = 2 \times 10^{-7} \times 2 \times
10^{-3}/10^{-2} = \qty{4e-8}{T}$ a \qty{10}{cm} --- poco, pero el coaxial da
cero.

\textbf{8.} $D = \sqrt{h^2 + d^2/4} = \qty{10.0}{m}$: cada uno da $\mu_0I/2\pi D =
\qty{1.0e-5}{T}$, perpendicular a la recta que va del conductor al punto. Al
ser opuestas las corrientes, las componentes paralelas a la recta que une
los conductores se cancelan y las verticales se suman: $B = 2 \times 10^{-5} \times
(d/2)/D = \qty{1.0e-6}{T}$, vertical. Cincuenta veces por debajo del campo terrestre
y cien veces por debajo de la recomendación.

\textbf{9.} Para $D \gg d$, $B \approx \dfrac{\mu_0I}{2\pi}\Bigl(\dfrac1{D_1} - \dfrac1{D_2}\Bigr)
\sim \dfrac{\mu_0I}{2\pi}\dfrac{d}{D^2}$: los dos campos en $1/D$ se cancelan a
primer orden y solo sobrevive su diferencia, de tamaño relativo $d/D$ --- un
``dipolo de dos hilos''.

\textbf{10.} El propio conductor a \qty{1}{m}: \qty{1.0e-4}{T} horizontal; el otro
a \qty{1.41}{m}: \qty{7.1e-5}{T} a \qty{45}{\degree}, en sentido opuesto: resultante
$(5 \times 10^{-5}, 5 \times 10^{-5})$, \qty{7e-5}{T} --- \qty{70}{\micro T}, por debajo de la
recomendación para el público y muy por debajo de la profesional.

\textbf{11.} $F/\ell = \mu_0I^2/2\pi d = 2 \times 10^{-7} \times 2.5 \times 10^5 =
\qty{0.05}{N/m}$, repulsiva; a \qty{20}{kA}: \qty{80}{N/m}.

\textbf{12.} \qty{15}{N} en régimen normal; \qty{24}{kN} en cortocircuito, ocho veces
el peso del vano, \qty{2.9}{kN}: los conductores salen despedidos, oscilan
y pueden chocar al volver --- de ahí los separadores y los interruptores que
abren en unos pocos ciclos.

\textbf{13.} $I\cos\omega t + I\cos(\omega t - 2\pi/3) + I\cos(\omega t + 2\pi/3) =
I\cos\omega t\,(1 + 2\cos2\pi/3) = 0$: los tres \omterm{def:b1:sinusoidal-impedance:complex}{fasores} suman cero. Los
campos en $1/D$ se cancelan como en el coaxial; el campo lejano decrece como $1/D^2$ o
más deprisa.

\textbf{14.} Allí el campo de la línea es vertical, y una brújula responde
solo a la componente horizontal: no hay desviación (la inclinación cambia en
$10^{-6}/(2 \times 10^{-5})$, tres grados, que la brújula no puede mostrar).

\textbf{15.} Lados $a$ (paralelos al eje): $\vect B$, radial, es
perpendicular a ellos; fuerza $NIaB$ en cada uno, perpendicular al plano
del cuadro y opuesta en los dos lados: un par. Lados $b$: $\dd\vect l$
es radial, paralelo a $\vect B$: sin fuerza.

\textbf{16.} $\Gamma = 2 \times NIaB \times b/2 = NIabB = NISB = 0.024\,I$ (N\,m, con $I$
en A). El campo radial está siempre en el plano del cuadro y es
perpendicular a los lados $a$: las fuerzas son siempre normales al plano y
su brazo $b/2$ no cambia nunca.

\textbf{17.} $C\theta = NISB$: $\theta = 2400\,I$ rad, es decir \qty{2.4}{rad/mA}; fondo
de escala \qty{1.2}{rad} con $I = \qty{0.50}{mA}$.

\textbf{18.} Longitud $200 \times 2(0.02 + 0.03) = \qty{20}{m}$, sección $\pi(5 \times
10^{-5})^2 = \qty{7.9e-9}{m^2}$: $R = 1.7 \times 10^{-8} \times 20/7.9 \times 10^{-9} =
\qty{43}{\ohm}$; $U = 43 \times 5 \times 10^{-4} = \qty{22}{mV}$; $P = UI = \qty{11}{\micro W}$.

\textbf{19.} $J\ddot\theta = -C\theta + NISB\,I$: $\omega_0 = \sqrt{C/J} = \sqrt{10^{-5}/6.7
\times 10^{-8}} = \qty{12}{rad/s}$, $T = \qty{0.51}{s}$.

\textbf{20.} $J\ddot\theta + \alpha\dot\theta + C\theta = NISB\,I$: crítico para
$\alpha = 2\sqrt{JC} = 2\sqrt{6.7 \times 10^{-13}} = \qty{1.6e-6}{N.m.s}$; la aguja
alcanza entonces su lectura lo más deprisa posible sin sobrepasarla.

\textbf{21.} Voltímetro: $R_s = 10/5 \times 10^{-4} - 43 \approx \qty{20}{k\ohm}$.
Amperímetro: la derivación lleva \qty{0.9995}{A} bajo \qty{22}{mV}: $R_{\text{sh}} =
\qty{22}{m\ohm}$.

\textbf{22.} $m = NIS = 200 \times 5 \times 10^{-4} \times 6 \times 10^{-4} = \qty{6e-5}{A.m^2}$;
$-mB = \qty{-1.2e-5}{J}$; resorte $\tfrac12C\theta^2 = \tfrac12 \times 10^{-5} \times 1.44 =
\qty{7.2e-6}{J}$: el mismo orden --- el resorte almacena lo que el par
magnético trabajó mientras el cuadro giraba.

\textbf{23.} En un campo uniforme el par vale $mB\cos\theta$ (máximo cuando
el plano contiene a $\vect B$ y nulo cuando la normal va según él): $C\theta =
NISB\,I\cos\theta$ --- lineal solo para $\theta$ pequeños y comprimida en lo
alto de la escala.

\textbf{24.} \qty{50}{Hz} frente a $\omega_0/2\pi = \qty{2}{Hz}$: el cuadro no puede
seguirlos; se queda en la media de la corriente, cero, con un temblor de
\omterm{def:b1:wave-propagation:signal}{amplitud} $\approx \theta_{\text{est}}(\omega_0/\omega)^2 = 1.2 \times (12/314)^2 \approx
\qty{2e-3}{rad}$, invisible. Un aparato de cuadro móvil lee la media --- para
la corriente alterna necesita un rectificador.

\textbf{25.} \qty{4e-4}{T} en el alma de un coaxial de \qty{1}{A} y nada fuera
(\omterm{thm:b1:magnetostatics:ampere}{teorema de Ampère}); \qty{1}{\micro T} bajo una línea de \qty{500}{A} y \qty{0.05}{N/m}
entre sus conductores (Biot--Savart más Laplace); \qty{2.4}{rad/mA} para
el aparato (par de Laplace contra un resorte).
\end{solution}
